%%%%%%%%%%%%Outline%%%%%%%%%%%%%%%%%%

%%%
% message<implementation is a small kernel core, a compile-time recognizer, and operation modules>
% we implement \tool in three parts
% the kernel core is a one-time change to the verifier and JIT
% the recognizer rewrites idioms in userspace before load
% each operation ships as one or more loadable kernel modules

%%% Section 6.1 Kernel Core

%%%
% message<the kernel core is a small one-time change>
% the kernel core adds 929 lines to the Linux kernel
% the table breaks the total down by component

%%%
% message<instruction encoding: a call plus a sidecar marked via src_reg>
% each extended instruction is a pair of instructions, a call and a sidecar
% both reuse existing eBPF opcodes, a reserved src_reg value marks the pair
% the call identifies the descriptor
% the sidecar carries a packed 52-bit payload with the operands
% the low four payload bits must form a valid eBPF register number
% the recognizer re-encodes payloads whose low bits exceed the largest register
% the kernel reverses the re-encoding when reading the sidecar

%%%
% message<verifier: a lowering stage and a restore stage around the existing analysis>
% the verifier gains a lowering stage and a restore stage around its analysis
% the lowering stage walks backward and substitutes each extended instruction's proof sequence
% the stage records each original pair for restoration
% the verifier then runs unchanged on the lowered program
% the restore stage reinstates the extended instructions after success
% with proof-sequence validation and bookkeeping the verifier changes total 487 lines

%%%
% message<JIT: per-architecture dispatch, fallback resolved before the JIT>
% the x86-64 and ARM64 JITs each gain a dispatch case for extended instructions
% the JIT looks up the descriptor and decodes the sidecar payload
% the JIT then calls the descriptor's native emit for the running architecture
% with no native emit the kernel rewrites the site to its proof sequence in load-time fixups
% dispatch is 97 lines on x86-64 and 121 on ARM64

%%%
% message<module registration reuses BTF and kfunc machinery>
% a loaded module registers its descriptors for lookup by the verifier and JIT
% registration reuses BTF, the kernel's type metadata, and kfuncs, kernel functions eBPF may call
% each module exports a stub kfunc as its BTF anchor
% the module lists its descriptors in a parallel array
% the kernel caches the descriptors per program
% BTF and registration are 139 lines
% the new descriptor type and helpers add 85 lines of headers

%%% Section 6.2 Compile-Time Recognizer

%%%
% message<the recognizer is an optional pre-load rewriter with two paths>
% the recognizer rewrites idioms into extended instructions before load, through two paths
% a program built without the recognizer, or unmatched, produces the same vanilla eBPF bytecode

%%%
% message<the selector path matches idioms during LLVM instruction selection>
% the first path is a selector in the LLVM eBPF backend
% the selector matches idioms during instruction selection and emits the call and sidecar
% counted over the backend's \tool commits, the selector adds 2,126 lines, half a new selection module

%%%
% message<the bytecode-recovery path re-matches idioms on compiled bytecode>
% the second path is a bytecode-recovery pass on already-compiled bytecode
% the pass covers bytecode the selector never sees or cannot match
% the implementation is a standalone userspace tool of 6,967 lines
% the pass checks that each rewrite preserves the program's semantics
% even a wrong rewrite leaves the program safe, the verifier re-checks the lowered program
% at worst a wrong rewrite changes the program's result

%%% Section 6.3 Operation Modules

%%%
% message<a module packages an operation's two forms outside the kernel core>
% each operation is packaged as one or more kernel modules, separate from the core
% a module supplies the proof-sequence routine, the native emits, and the BTF registration

%%%
% message<modules are organized by native instruction; tree counts and scope>
% each module covers one native instruction on one architecture
% a single-architecture operation contributes modules on one side, so the counts differ
% the tree is 14 x86-64 and 11 ARM64 modules, 9,765 lines, including helpers and test-only modules

%%%%%%%%%%%%Outline%%%%%%%%%%%%%%%%%%

\section{Implementation}
\label{sec:implementation}
% 实现

We implement \tool in three parts.
% 我们分三部分实现 \tool。
The kernel core is a one-time change to the verifier and the JIT (\S\ref{subsec:kernel}).
% 内核核心是对验证器和 JIT 的一次性更改。
A compile-time recognizer rewrites idioms in userspace before load (\S\ref{subsec:llvm}).
% 编译时识别器在加载前在用户空间重写习语。
Each operation ships as one or more loadable kernel modules (\S\ref{subsec:descriptors}).
% 每个操作作为一个或多个可加载内核模块发布。

\subsection{Kernel Core}
\label{subsec:kernel}
% 内核核心

The kernel core adds 929 lines to the Linux kernel.
% 内核核心向 Linux 内核添加了 929 行。
Table~\ref{tab:kernel-loc} breaks the total down by component.
% 表按组件分解总数。

\begin{table}[t]
\centering
\small
\caption{The kernel core is a 929-line addition, measured as lines added on the \tool kernel branch over mainline Linux. The total excludes 11 lines of disassembler support and tools-side header sync.}
\label{tab:kernel-loc}
\begin{tabular}{@{}lc@{}}
\toprule
Component & Lines of code \\
\midrule
Verifier (lowering, restore, validation) & 487 \\
BTF and registration & 139 \\
JIT dispatch (ARM64) & 121 \\
JIT dispatch (x86-64) & 97 \\
Headers & 85 \\
\midrule
Total kernel core & 929 \\
\bottomrule
\end{tabular}
\end{table}

\noindent\textbf{Instruction encoding.}
% \noindent\textbf{指令编码。}
Each extended instruction is a pair of instructions, a call and a sidecar, shown in Table~\ref{tab:encoding}.
% 每条扩展指令是一对指令：一个调用和一个伴随，如表所示。
Both reuse existing eBPF opcodes, and a reserved \texttt{src\_reg} value marks the pair for \tool.
% 两者都重用现有的 eBPF 操作码，保留的 src\_reg 值标记该对用于 \tool。
The call identifies the descriptor to use.
% 调用标识要使用的描述符。
The sidecar carries a packed 52-bit payload with the operands.
% 伴随携带包含操作数的 52 位压缩载荷。
The low four bits of the payload must form a valid eBPF register number.
% 载荷的低四位必须形成有效的 eBPF 寄存器编号。
The recognizer therefore re-encodes a payload whose low four bits exceed the largest register number.
% 因此识别器重新编码低四位超过最大寄存器编号的载荷。
The kernel reverses the re-encoding when reading the sidecar.
% 内核在读取伴随时逆转重新编码。

\begin{table}[t]
\centering
\small
\caption{Extended-instruction encoding. A reserved \texttt{src\_reg} value marks the pair for \tool, a \texttt{BPF\_CALL} for the call and a \texttt{MOV64\_K} for the sidecar. The notation [h:l] gives a field's position within the 52-bit sidecar payload.}
\label{tab:encoding}
\begin{tabular}{@{}llcp{0.50\columnwidth}@{}}
\toprule
Insn & Field & Bits & Meaning \tabularnewline
\midrule
Call & \texttt{imm} & 32 & \raggedright BTF ID of the operation's stub kfunc \tabularnewline
 & \texttt{off} & 16 & \raggedright Index into \texttt{fd\_array} for the module's BTF file descriptor \tabularnewline
\midrule
Sidecar & \texttt{dst\_reg} & [3:0] & \raggedright Destination register or per-descriptor form tag \tabularnewline
 & \texttt{off} & [19:4] & \raggedright Offset or auxiliary field per descriptor \tabularnewline
 & \texttt{imm} & [51:20] & \raggedright Immediate value or configuration per descriptor \tabularnewline
\bottomrule
\end{tabular}
\end{table}

\noindent\textbf{Verifier lowering and restore.}
% \noindent\textbf{验证器降低和恢复。}
The verifier gains a lowering stage and a restore stage around its existing analysis.
% 验证器在其现有分析周围获得降低阶段和恢复阶段。
The lowering stage traverses the program backward and replaces each extended instruction with the proof sequence from its descriptor.
% 降低阶段向后遍历程序，用描述符中的证明序列替换每条扩展指令。
The stage records each original pair for restoration.
% 该阶段记录每个原始对以供恢复。
The verifier then runs unchanged on the lowered program.
% 然后验证器在降低后的程序上不变地运行。
The restore stage reinstates the extended instructions after verification succeeds.
% 验证成功后恢复阶段恢复扩展指令。
With proof-sequence validation and the supporting bookkeeping, the changes to the verifier total 487 lines, the largest single component.
% 加上证明序列验证和支持簿记，验证器的更改共 487 行，是最大的单一组件。

\noindent\textbf{JIT dispatch.}
% \noindent\textbf{JIT 分派。}
The x86-64 and ARM64 JIT backends each gain a dispatch case for extended instructions.
% x86-64 和 ARM64 JIT 后端各获得扩展指令的分派情况。
On reaching an extended instruction, the JIT looks up the descriptor and decodes the sidecar payload.
% 到达扩展指令时，JIT 查找描述符并解码伴随载荷。
The JIT then calls the descriptor's native emit for the running architecture.
% 然后 JIT 调用描述符对应运行架构的本地发射。
When a descriptor has no native emit for the architecture, the kernel rewrites the site back to its proof sequence during load-time fixups, before the JIT runs.
% 当描述符没有该架构的本地发射时，内核在 JIT 运行前的加载时修复期间将站点重写回其证明序列。
The dispatch logic is 97 lines on x86-64 and 121 lines on ARM64.
% 分派逻辑在 x86-64 上 97 行，在 ARM64 上 121 行。

\noindent\textbf{Module registration.}
% \noindent\textbf{模块注册。}
A loaded module registers its descriptors for lookup by the verifier and the JIT.
% 加载的模块注册其描述符以供验证器和 JIT 查找。
Registration reuses two existing kernel facilities, BTF, the kernel's type metadata for eBPF objects, and kfuncs, kernel functions that eBPF programs may call.
% 注册重用两个现有内核设施：BTF（内核对 eBPF 对象的类型元数据）和 kfunc（eBPF 程序可以调用的内核函数）。
Each module exports a stub kfunc as its BTF anchor.
% 每个模块导出一个存根 kfunc 作为其 BTF 锚点。
The module lists its \texttt{struct bpf\_kinsn} descriptors in a parallel array.
% 模块在并行数组中列出其 struct bpf\_kinsn 描述符。
The kernel caches the descriptors per program.
% 内核为每个程序缓存描述符。
The BTF and registration changes are 139 lines.
% BTF 和注册更改为 139 行。
The new \texttt{struct bpf\_kinsn} type and helpers add 85 lines of headers.
% 新的 struct bpf\_kinsn 类型和辅助函数添加 85 行头文件。

\subsection{Compile-Time Recognizer}
\label{subsec:llvm}
% 编译时识别器

The recognizer rewrites a program's idioms into extended instructions before load, through two paths.
% 识别器通过两条路径在加载前将程序的习语重写为扩展指令。
A program built without the recognizer, or one whose idioms the recognizer does not match, produces the same vanilla eBPF bytecode and runs through the unchanged pipeline.
% 没有使用识别器构建的程序，或识别器不匹配其习语的程序，产生相同的原生 eBPF 字节码并通过不变的管线运行。

\noindent\textbf{Selector.}
% \noindent\textbf{选择器。}
The first path is a selector in the LLVM eBPF backend.
% 第一条路径是 LLVM eBPF 后端中的选择器。
The selector matches idioms during instruction selection and emits the call and sidecar in place of the decomposed bytecode.
% 选择器在指令选择期间匹配习语，并发射调用和伴随替代分解的字节码。
Counted over the backend's \tool commits, the selector adds 2{,}126 lines, half in a new selection module.
% 计算后端的 \tool 提交，选择器添加 2126 行，一半在新的选择模块中。

\noindent\textbf{Bytecode recovery.}
% \noindent\textbf{字节码恢复。}
The second path is a bytecode-recovery pass that re-matches idioms on already-compiled bytecode.
% 第二条路径是字节码恢复遍，在已编译的字节码上重新匹配习语。
The pass covers bytecode the selector never sees or cannot match, such as programs compiled by a separate toolchain.
% 该遍覆盖选择器从未见过或无法匹配的字节码，如由单独工具链编译的程序。
The implementation is a standalone userspace tool of 6{,}967 lines.
% 实现是一个 6967 行的独立用户空间工具。
The pass checks that each rewrite preserves the program's semantics, for example that a temporary register the idiom overwrites is dead afterward.
% 该遍检查每次重写保持程序的语义，例如习语覆盖的临时寄存器之后是死的。
Even a wrong rewrite leaves the program safe, because the verifier re-checks every lowered program (\S\ref{subsec:safety}).% At worst, a wrong rewrite changes the program's result.
% 即使错误的重写也使程序保持安全，因为验证器重新检查每个降低的程序。

\subsection{Operation Modules}
\label{subsec:descriptors}
% 操作模块

Each operation is packaged as one or more kernel modules, separate from the kernel core.
% 每个操作打包为一个或多个内核模块，与内核核心分开。
A module supplies an operation's proof-sequence routine, its native emits, and the BTF registration that exposes both.
% 模块提供操作的证明序列例程、其本地发射和暴露两者的 BTF 注册。

Each module covers one native instruction on one architecture.
% 每个模块覆盖一个架构上的一条本地指令。
An operation that targets a single architecture contributes modules on that architecture only, which is why the per-architecture counts differ.
% 针对单一架构的操作只在该架构上贡献模块，这就是每架构计数不同的原因。
The module tree holds 14 modules on x86-64 and 11 on ARM64, 9{,}765 lines in all, including shared helper headers, test-only modules, and operations beyond \lang's seven (Table~\ref{tab:descriptors}).
% 模块树在 x86-64 上有 14 个模块，在 ARM64 上有 11 个，共 9765 行，包括共享辅助头文件、仅测试模块和超出 \lang 七个的操作。
